% GENERATED FILE. Edit canonical lessons, then run npm run generate:books.
% canonical-source-hash: fnv1a64:2d5bbd65be953f54
% canonical-lessons: FR-C367-supposer, FR-C367-hypothese, FR-C367-presumer, FR-C367-plausible, FR-C367-supposition

\chapter{Supposons}
\label{ch:fr-supposons}
\hlchaptermodality{367}

\begin{chapteropening}
\textbf{By the end of this chapter:} \emph{I can make a hypothesis: suppose or presume something, and say an explanation is plausible.}

Complete the last lesson of chapter 367: I can make a hypothesis: suppose or presume something, and say an explanation is plausible.
\end{chapteropening}

\section[supposer]{supposer — to suppose, to assume}

\label{lesson:FR-C367-supposer}



\begin{quote}
Before the new one: say the French for the solidarity, then the French for the community.
\end{quote}

\subsection*{You'll want to know: supposer}
\textbf{supposer} — \textquotedblleft{}to suppose, to assume\textquotedblright{}.

\begin{quote}
Je suppose que tu es fatigué après ce voyage.
\end{quote}

\begin{quote}
\emph{I suppose you are tired after that trip.}
\end{quote}

\emph{Supposer que} + indicative. \emph{À supposer que} takes the subjunctive.

\begin{cousinweb}
\emph{supposer} comes through Middle French (EtymDB) from Latin \emph{suppōnere} \textquotedblleft{}to put under, to put in place of\textquotedblright{} (Lewis \& Short), which took the sense \textquotedblleft{}to assume\textquotedblright{} in Medieval Latin; it is built on \emph{pōnere} \textquotedblleft{}to put\textquotedblright{}.
\end{cousinweb}

\begin{grammarlens}[title={What you've built}]
One more word for telling, explaining and describing at B1.
\end{grammarlens}

\subsection*{Your turn}
\begin{itemize}
\raggedright
  \item \emph{Say it:} \emph{supposer}
  \item \emph{Say it:} \emph{supposer}, once more
  \item \emph{Say it:} a sentence of your own with \emph{supposer}
  \item \emph{Recall:} say the French for the solidarity, then the French for the community, then say \emph{supposer} again
\end{itemize}

\begin{tcolorbox}[breakable,colback=teal!4,colframe=teal!35!black,title={Before you move on}]
What does supposer mean? (\textquotedblleft{}To suppose, to assume\textquotedblright{}.) Say it once more.
\end{tcolorbox}
\section[l'hypothèse]{l'hypothèse — the hypothesis}

\label{lesson:FR-C367-hypothese}



\begin{quote}
Before the new one: say the French for the community, then the French for to suppose.
\end{quote}

\subsection*{You'll want to know: l'hypothèse}
\textbf{l'hypothèse} — \textquotedblleft{}the hypothesis\textquotedblright{}.

\begin{quote}
Dans cette hypothèse, nous partirions plus tôt.
\end{quote}

\begin{quote}
\emph{On that hypothesis, we would leave earlier.}
\end{quote}

\emph{Dans l'hypothèse où} takes the conditional. Feminine.

\begin{cousinweb}
\emph{hypothèse} is a learned borrowing from Latin \emph{hypothesis} (EtymDB), from Greek \emph{hupóthesis} \textquotedblleft{}a foundation, a supposition\textquotedblright{}, built from \emph{hupó} \textquotedblleft{}under\textquotedblright{} and \emph{thésis} \textquotedblleft{}a placing\textquotedblright{}.
\end{cousinweb}

\begin{grammarlens}[title={What you've built}]
One more word for telling, explaining and describing at B1.
\end{grammarlens}

\subsection*{Your turn}
\begin{itemize}
\raggedright
  \item \emph{Say it:} \emph{l'hypothèse}
  \item \emph{Say it:} \emph{l'hypothèse}, once more
  \item \emph{Say it:} a sentence of your own with \emph{l'hypothèse}
  \item \emph{Recall:} say the French for the community, then the French for to suppose, then say \emph{l'hypothèse} again
\end{itemize}

\begin{tcolorbox}[breakable,colback=teal!4,colframe=teal!35!black,title={Before you move on}]
What does l'hypothèse mean? (\textquotedblleft{}The hypothesis\textquotedblright{}.) Say it once more.
\end{tcolorbox}
\section[présumer]{présumer — to presume}

\label{lesson:FR-C367-presumer}



\begin{quote}
Before the new one: say the French for to suppose, then the French for the hypothesis.
\end{quote}

\subsection*{You'll want to know: présumer}
\textbf{présumer} — \textquotedblleft{}to presume\textquotedblright{}.

\begin{quote}
Je présume que le magasin est fermé le dimanche.
\end{quote}

\begin{quote}
\emph{I presume the shop is closed on Sundays.}
\end{quote}

Formal: \emph{je présume que oui} — I presume so.

\begin{cousinweb}
\emph{présumer} is a learned borrowing from Latin \emph{praesūmere} \textquotedblleft{}to take beforehand, to suppose\textquotedblright{} (Lewis \& Short), built on \emph{sūmere} \textquotedblleft{}to take\textquotedblright{}.
\end{cousinweb}

\begin{grammarlens}[title={What you've built}]
One more word for telling, explaining and describing at B1.
\end{grammarlens}

\subsection*{Your turn}
\begin{itemize}
\raggedright
  \item \emph{Say it:} \emph{présumer}
  \item \emph{Say it:} \emph{présumer}, once more
  \item \emph{Say it:} a sentence of your own with \emph{présumer}
  \item \emph{Recall:} say the French for to suppose, then the French for the hypothesis, then say \emph{présumer} again
\end{itemize}

\begin{tcolorbox}[breakable,colback=teal!4,colframe=teal!35!black,title={Before you move on}]
What does présumer mean? (\textquotedblleft{}To presume\textquotedblright{}.) Say it once more.
\end{tcolorbox}
\section[plausible]{plausible — plausible}

\label{lesson:FR-C367-plausible}



\begin{quote}
Before the new one: say the French for the hypothesis, then the French for to presume.
\end{quote}

\subsection*{You'll want to know: plausible}
\textbf{plausible} — \textquotedblleft{}plausible\textquotedblright{}.

\begin{quote}
Son explication est plausible.
\end{quote}

\begin{quote}
\emph{His explanation is plausible.}
\end{quote}

One form for both genders.

\begin{cousinweb}
\emph{plausible} is a learned borrowing from Latin \emph{plausibilis} \textquotedblleft{}deserving applause, acceptable\textquotedblright{} (Lewis \& Short), from \emph{plaudere} \textquotedblleft{}to clap\textquotedblright{}.
\end{cousinweb}

\begin{grammarlens}[title={What you've built}]
One more word for telling, explaining and describing at B1.
\end{grammarlens}

\subsection*{Your turn}
\begin{itemize}
\raggedright
  \item \emph{Say it:} \emph{plausible}
  \item \emph{Say it:} \emph{plausible}, once more
  \item \emph{Say it:} a sentence of your own with \emph{plausible}
  \item \emph{Recall:} say the French for the hypothesis, then the French for to presume, then say \emph{plausible} again
\end{itemize}

\begin{tcolorbox}[breakable,colback=teal!4,colframe=teal!35!black,title={Before you move on}]
What does plausible mean? (\textquotedblleft{}Plausible\textquotedblright{}.) Say it once more.
\end{tcolorbox}
\section[la supposition]{la supposition — the supposition, the assumption}

\label{lesson:FR-C367-supposition}



\begin{quote}
Before the new one: say the French for to presume, then the French for plausible.
\end{quote}

\subsection*{You'll want to know: la supposition}
\textbf{la supposition} — \textquotedblleft{}the supposition, the assumption\textquotedblright{}.

\begin{quote}
Ce n'est qu'une supposition : je n'en suis pas sûr.
\end{quote}

\begin{quote}
\emph{It is only an assumption: I am not sure of it.}
\end{quote}

\emph{Faire une supposition} — to make an assumption.

\begin{cousinweb}
\emph{supposition} is a learned borrowing from Latin \emph{suppositiō} \textquotedblleft{}a putting under, a substitution\textquotedblright{} (Lewis \& Short), from \emph{suppōnere}, built on \emph{pōnere} \textquotedblleft{}to put\textquotedblright{}.
\end{cousinweb}

\begin{grammarlens}[title={What you've built}]
One more word for telling, explaining and describing at B1.
\end{grammarlens}

\subsection*{Your turn}
\begin{itemize}
\raggedright
  \item \emph{Say it:} \emph{la supposition}
  \item \emph{Say it:} \emph{la supposition}, once more
  \item \emph{Say it:} a sentence of your own with \emph{la supposition}
  \item \emph{Recall:} say the French for to presume, then the French for plausible, then say \emph{la supposition} again
\end{itemize}

\begin{tcolorbox}[breakable,colback=teal!4,colframe=teal!35!black,title={Before you move on}]
What does la supposition mean? (\textquotedblleft{}The supposition, the assumption\textquotedblright{}.) Say it once more.
\end{tcolorbox}
